\section{Background and Related Work}
\label{sec:background}

We synthesize four strands of literature that bracket the
contribution: DeFi lending mechanics, the HFT microstructure canon
we borrow from, MCDM in algorithmic allocation, and the equilibrium
impossibility result of \cite{halpern2024fair} as the most relevant
theoretical counter-evidence.

\subsection{DeFi lending mechanics}
\label{sec:bg-defi-mechanics}

Aave V3, Compound V3, Morpho Blue, Euler V2, Spark, and Fluid all
implement deterministic interest-rate models in which the supply
APY is a piecewise-linear (or in Morpho Blue's case, adaptive)
function $f_{\text{kink}}(u)$ of pool utilization
$u = \text{borrowed} / \text{supplied}$.  The function has a
shallow slope below a kink point (typically $u^{*} = 0.8$--$0.9$)
and a steep slope above it that taxes high utilization.  Although
$f_{\text{kink}}$ is deterministic given $u$, the realised rate
stream is stochastic because $u$ moves with on-chain demand at
per-block granularity --- exactly the resolution at which our
allocator operates.

The empirical structure of the cross-protocol spread is critical.
\cite{gudgeon2020loanable} establish cointegration between Aave
and Compound USDC rates with Compound leading and Aave adjusting
at speed $0.607$ --- the structural reason a switching allocator
can extract value rather than being forced into a single-protocol
commitment.  Recent cross-chain comparative analyses
\cite{aavecompoundcrosschain2025} confirm the cointegration
persists into 2026.  On our \NBlocksFull-block panel
(\Cref{sec:empirical}) the spread structure crosses regime around
2025-Q4 with realised cross-protocol volatility jumping by a
factor of five between Q1 and Q2 of 2026 --- an adversarial split
that tests transferability under distribution shift.

\subsubsection{Market structure (April 2026 snapshot)}
\label{sec:bg-market-structure}

The DeFi lending market is highly concentrated.  As of April 2026,
total deposits across all lending protocols sum to roughly
\$54\,B, with the top ten protocols accounting for $\sim 78\%$ of
that value \cite{defillama2026snapshot}.  \Cref{tab:lending-tvl}
reports the top-eight breakdown; the three protocols highlighted
in bold delimit the scope of this paper's empirical study, jointly
holding $\sim 47\%$ of total lending TVL.

\begin{table}[t]
  \centering
  \caption{Top-eight DeFi lending protocols by total value locked
    (TVL), April 2026 \cite{defillama2026snapshot}.  Bold rows mark
    the subset covered by this paper.}
  \label{tab:lending-tvl}
  \begin{tabular}{lrr}
    \toprule
    Protocol (chain)            & TVL (\$B) & Share \\
    \midrule
    \textbf{Aave V3}            & \textbf{19.4} & \textbf{36.0\%} \\
    Spark (SparkLend)           & 6.8           & 12.6\% \\
    \textbf{Morpho Blue}        & \textbf{4.9}  & \textbf{9.1\%}  \\
    Compound V3                 & 2.7           & 5.0\% \\
    JustLend (Tron)             & 2.4           & 4.4\%  \\
    Fluid                       & 1.6           & 3.0\%  \\
    Kamino (Solana)             & 1.1           & 2.0\%  \\
    \textbf{Euler V2}           & \textbf{0.89} & \textbf{1.6\%}  \\
    \midrule
    Top-8 subtotal              & 39.8          & 73.7\% \\
    Tail ($\sim 350+$ protocols)& 14.2          & 26.3\% \\
    \bottomrule
  \end{tabular}
\end{table}

\subsection{HFT microstructure analogs}
\label{sec:bg-hft}

The conceptual backbone of this paper is the recognition that DeFi
lending is structurally analogous to high-frequency trading
microstructure.  Five source books anchor the methodology, each
mapped to one section of this paper.

\cite{ohara1995} provides the adverse-selection framework (Glosten-
Milgrom, Kyle): the Kyle-$\lambda$ market impact parameter is
literally $\partial r / \partial u$ on the IRM curve in DeFi ---
\emph{observable}, not inferred --- and the Demsetz price of
immediacy inverts (depositors paid by borrowers, not the reverse).
\cite{krause2005} contributes the closed-form market-depth
expression $1/\lambda = TVL \cdot (1 - u) / \text{slope}_{1}$ and
the resilience half-life decomposition that calibrates T2's OU
mean-reversion rate $\kappa$.  \cite{kissell2014} contributes the
Implementation Shortfall framework adapted as
$\textsc{IS}_{\text{defi}} = \int (r^{*} - r_{\text{held}}) V \, dt
+ \sum (\text{gas} + \text{slip} + \text{mev})$ and the
closed-form optimal trade rate (eq.\,8.23, p.\,281) that benchmarks
T2.  \cite{lopezdeprado2018} contributes the triple-barrier method
that maps directly to switch/hold/timeout, the purged k-fold with
embargo that we use for T3's cross-validation, and the Deflated
Sharpe Ratio threshold ($\text{DSR} > 0.95$) that gates our
multi-hypothesis tests.  \cite{mackenzie2021} contributes the
four-class signal taxonomy (Table 3.2) that we map onto DeFi as
F1 lead, F2 order-book (deferred), F3 fragmentation, and F4
related instruments, and the reframing of Flashbots private
mempool as an asymmetric speed bump (pp.\,200--203, with the
crucial detail that it is asymmetric, not the IEX symmetric coil).

\subsection{Algorithmic allocation in DeFi rates}
\label{sec:bg-forecasting}

Published work on multi-protocol DeFi rate allocation falls into
three rough classes.  \emph{Protocol-side rate control}:
\cite{bertucci2025rl} train offline RL agents (CQL, BC, TD3-BC) on
historical Aave data to set rates from the protocol governance
side; the closely related Auto.gov framework \cite{xu2023autogov}
reports $\sim 14\%$ improvement over benchmarks under RL.
\emph{Adaptive rate control}: \cite{agilerate2024} fits recursive
least-squares controllers on Compound demand/supply curves at
3-hour resolution.  \emph{Optimal protocol design}:
\cite{baude2025optimal} derives a closed-form Riccati-ODE solution
for the protocol-optimal rate curve under linear agent response
and a Monte-Carlo + deep-learning estimator for the nonlinear case;
block-level empirical calibration shows industry-standard
piecewise rate curves are sub-optimal, leaving structural
residuals.  This last point is foundational for us:
\cite{baude2025optimal} solves the \emph{protocol designer's}
problem (which rate function should the contract use?), whereas we
solve the \emph{lender's} problem (given the realised rate stream,
where should liquidity go?).  We are not aware of a published
event-time multi-protocol allocator in the May 2026 literature;
the nearest neighbour is the LSTM+Q-learning AMM quoter of
\cite{predamm2024}, which transplants a TradFi forecaster to
Uniswap V3 but does not span multiple lending protocols.

\subsection{MCDM in algorithmic allocation}
\label{sec:bg-mcdm}

Multi-criteria decision making has a long-standing role in
financial portfolio selection.  TOPSIS \cite{hwang1981mcdm} and
PROMETHEE \cite{brans1985promethee} are the two dominant families.
Recent cryptocurrency-specific applications include
\cite{aljinovic2021crypto}'s multicriteria crypto portfolio
selection and \cite{hosseinzadeh2023promethee}'s asymmetric
PROMETHEE-II that integrates return prediction with multi-criteria
ranking.  Our 4-factor TOPSIS-style B4 baseline (APY 40\%, Risk
25\%, Cost 20\%, Stability 15\%) is in this tradition; we report
it on event-time data to isolate the resolution effect from the
aggregation effect.

\subsection{Counter-evidence: Halpern--Pass--Saraf impossibility}
\label{sec:bg-impossibility}

The most pointed theoretical counter-evidence to active DeFi
lending strategies is \cite{halpern2024fair}.  Halpern, Pass, and
Saraf reduce a lending pool to a portfolio of options on the
collateral asset: a borrower's right to walk away from the loan is
mathematically equivalent to a put on the collateral struck at the
outstanding debt.  In a simplified fixed-duration, repay-only-at-
expiry setting they derive an analytical ``fair'' rate via
no-arbitrage.  They then prove that in the realistic setting ---
dynamic, perpetual loans with early repayment of the kind Aave V3
and Compound V3 implement --- \emph{no fair interest rate exists}:
any rate either over-compensates lenders relative to the
option-implied premium or under-compensates them, because the
borrower's repayment optionality has no finite hedging cost.

A naive reading of \cite{halpern2024fair} suggests there is no
exploitable structure in DeFi rates and any allocator should
underperform a passive benchmark.  We argue the implication does
not follow.  Halpern et al.\ prove a \emph{static equilibrium
fairness} impossibility under no-arbitrage --- a statement about
whether a single rate function can correctly price the borrower-
side optionality in equilibrium.  Our empirical claim is the
orthogonal \emph{event-time predictability} statement: the
realised rate stream observed off-chain at per-block resolution
exhibits enough short-horizon structure (cointegration, regime
persistence, kink-residual cycles) to drive a gas-aware switching
strategy that significantly outperforms passive single-protocol
hold under a paired monthly Sharpe bootstrap.  Equivalently, the
absence of a single equilibrium-fair rate is consistent with ---
and indeed implies --- the persistent disequilibrium that an
event-time allocator can target.  The two claims constrain
different layers of the protocol stack: \cite{halpern2024fair}
constrains pricing theory at the contract design layer; we operate
at the lender-side execution layer.
